\section{Division-free quadratic arithmetic and bit complexity}
\label{sec:exact}

\subsection{An exact ring for each integer color count}

Fix an integer $q\ge5$ and put
\[
 R_q=\Z[x]/(x^2-(q-2)x+1).
\]
The discriminant $(q-2)^2-4$ lies strictly between the consecutive squares
$(q-3)^2$ and $(q-2)^2$, so the polynomial is irreducible over $\Q$.
In either real embedding,
\[
 x+x^{-1}=q-2,\qquad x^{-1}=q-2-x,\qquad (x+1)^2=qx.
\]
In particular $x$ is a unit. Each element has a unique representation
$a+bx$ with $a,b\in\Z$. Pair inequality is exact inequality of the algebraic
numbers; it is not a numerical tolerance test.

For $q=6$, the positive roots are $2+\sqrt3$ and $2-\sqrt3$.
The Jones variable in our derivation is $t=A^{-4}=x^{-1}$.
Changing embeddings interchanges these reciprocal values; equality to $1$
is preserved because the coefficients are rational integers.
The code works in the ring and need not select a real embedding.

\begin{table}[htbp]
\centering\small
\begin{tabular}{@{}ll@{}}
\toprule
Operation & Pair formula\\
\midrule
Addition & $(a,b)+(c,d)=(a+c,b+d)$\\
Multiplication & $(a,b)(c,d)=(ac-bd,\ ad+bc+(q-2)bd)$\\
Multiply by $-x$ & $(a,b)\mapsto(b,\ -a-(q-2)b)$\\
Multiply by $-x^{-1}$ & $(a,b)\mapsto(-(q-2)a-b,\ a)$\\
Inverse of $x$ & $x^{-1}=(q-2,-1)$\\
\bottomrule
\end{tabular}
\caption{Exact arithmetic used by the transfer. Full pair multiplication
is needed for normalization and component merges; a single-edge transfer
uses only additions and small integer multiples.}
\label{tab:pairs}
\end{table}

Equations~\eqref{eq:interaction} and \eqref{eq:unknot-partition} now define a
division-free one-tensor algorithm. Negative powers in normalization are
powers of the explicit unit $x^{-1}$. Fast powering is exact. No prime choice,
probability-of-error estimate, or floating-point precision policy is required.

\begin{proposition}[Sound exact obstruction]\label{prop:obstruction}
If the partition function and $Z_U$ are unequal as pairs in $R_q$, then
$D$ is a nontrivial knot. If they are equal, the computation gives no
unknot certificate.
\end{proposition}
\begin{proof}
The nonzero normalization factor in \eqref{eq:normalize-E} identifies
partition equality with equality of the selected normalized Jones value to
$1$. Every unknot has normalized Jones polynomial $1$, so inequality rules
it out. The converse would require an additional detection theorem that
is not available here. Section~\ref{sec:blind} gives explicit failures at
$q=5$.
\end{proof}

\subsection{Coefficient lengths}

Let $\|a+bx\|_1=|a|+|b|$. The two equal-spin multiplication maps in
Table~\ref{tab:pairs} have induced $\ell^1$ norm at most $q-1$.

\begin{theorem}[Linear coefficient bit lengths]\label{thm:bits}
After $i$ edges and $I$ introduced vertices, every aggregate table
coefficient, including every partial sum inserted into the next table,
satisfies
\[
 \|a+bx\|_1\le q^I(q-1)^i\le q^v(q-1)^n.
\]
Its magnitude coefficients therefore have at most
\[
 \left\lfloor v\log_2q+n\log_2(q-1)\right\rfloor+1
\]
bits. For fixed $q$ this is $O(n)$, and for $q=6$ the simpler upper bound
$5n+4$ is valid.
\end{theorem}
\begin{proof}
Every labeled assignment of the introduced vertices contributes a product
of $i$ edge weights. Its norm is at most $(q-1)^i$.
An aggregate coefficient is a sum over a subset of at most $q^I$
assignments, so the triangle inequality proves the first bound. Partial
insertions are sums over subsets of the same histories and obey the same
bound. The connected Tait graph has $v\le n+1$, giving the stated bit length.
For $q=6$, $\log_2 6+\log_2 5=\log_2 30<5$ and
$\log_2 6<3$, proving the convenient integer bound.
\end{proof}

The normalization pair is also polynomial in bit length:
\[
 \|Z_U\|_1\le(q-1)^{|h|}q^{v+1},\qquad
 |h|\le n+(v+1)/2.
\]
Thus no hidden exponential coefficient size invalidates the transfer's
width bound. At fixed $q$, a single-edge update costs $O(n)$ elementary bit
operations apart from key management. Ordinary multiplication of
$O(n)$-bit integers suffices for powering and component merges; faster
integer multiplication can be used without changing the width dependence.

The coefficient-length bound is a worst-case upper bound, not an assertion
that every table entry reaches it. The experiment record reports the
largest coefficient actually inserted. It does not treat Python object
overhead as coefficient bits, and neither number is a process-memory limit.

\subsection{The modular comparison backend}

For a second implementation path the package includes a modular $q=5$
backend over $p=2^{61}-1$, with a root of $x^2-3x+1$ and a compatible fourth
root $A$. It compares the full bracket at precisely that $A$ with
$\delta(-A^3)^{\wr(D)}$. The constants are checked algebraically on import.
The old matching transfer can be evaluated at the same $A$, making an
apples-to-apples scalar comparison possible.

A modular mismatch is also a sound obstruction. A modular match may arise
from reduction modulo $p$ or from a genuine specialization collision.
Increasing the number of primes can address the former but cannot repair
the latter. In particular it cannot repair the family in
Section~\ref{sec:blind}. The exact backend makes this distinction observable.
